% Source: OpenAI2026PolynomialPEPS, Proposition 4.1, 03-patches.tex,
% eq:patch-variational-problem, eq:patch-clipped-eigenvalues, eq:patch-log-ratio;
% revision adc7f1241b42e322a6451854ab7e4b4c146bf78a.
\begin{lemma}[Clipped inverse-power filters at a minimum]
    \label{lem:peps_actual_patch_minimum_filter_clipping}
    \lean{TNLean.PEPS.exists_clipped_filter_of_isMinOn_regularizedPatchObjective}
    \uses{def:peps_regularized_patch_minimum,def:peps_patch_canonical_marginal}
    Let $V$ be finite, with local spaces $\mathcal H_v=\C^{n_v}$, and let
    $X_0\subseteq\cdots\subseteq X_{m-1}\subseteq V$.
    Suppose that $a_j>0$, $b>0$, and $\|\Omega\|=1$.
    For independent density matrices $x_j$ on $\mathcal H_{X_j}$, set
    \begin{align}
         & K_j=(x_j+bI_{X_j})^{-a_j/2},\qquad
        \Psi(x)=\widehat K_{m-1}\cdots\widehat K_0\Omega, \notag
    \end{align}
    where each hat denotes extension by the identity outside $X_j$.
    Suppose that $x$ minimizes $\|\Psi(x)\|$ over these density matrices,
    and write
    \begin{align}
         & \phi=\frac{\Psi(x)}{\|\Psi(x)\|},\qquad
        \rho_j=\operatorname{Tr}_{V\setminus X_j}|\phi\rangle\langle\phi|. \notag
    \end{align}
    Then $K_j\rho_j=\rho_jK_j$ for every $j$. There are a unitary $W_j$
    and real families $\ell_{j,i},p_{j,i}$ such that $\ell_{j,i}>0$ and
    \begin{align}
         & K_j=W_j\operatorname{diag}(\ell_{j,i})W_j^*,\qquad
        \rho_j=W_j\operatorname{diag}(p_{j,i})W_j^*. \notag
    \end{align}
    Whenever $p_{j,i},p_{j,k}>0$,
    \begin{align}
         & |\log\ell_{j,i}-\log\ell_{j,k}|
        \leq\frac{a_j}{2}|\log p_{j,i}-\log p_{j,k}|. \notag
    \end{align}
    Regions may be empty or repeated. This is the logarithmic filter
    estimate in \cite[Proposition~4.1]{OpenAI2026PolynomialPEPS}.
\end{lemma}
\begin{proof}
    \uses{thm:peps_ordered_patch_marginal_coordinates}
    At the minimum, $x_j$ and $\rho_j$ have a common unitary eigenbasis
    with eigenvalues $s_i,p_i$ satisfying
    $s_i+b=\max\{p_i/\lambda,b\}$ for some $\lambda>0$.
    Functional calculus gives positive filter eigenvalues
    $\ell_i=(s_i+b)^{-a_j/2}$ in the same basis, and hence commutation.
    The function $t\mapsto\max\{t,\log b\}$ is $1$-Lipschitz.
    Applying this to $\log(p_i/\lambda)$ and $\log(p_k/\lambda)$ gives
    the displayed bound. The coordinate identity for normalized marginals
    identifies the marginal in this argument with the stated $\rho_j$.
\end{proof}
